\documentclass[11pt,a4paper]{article}
\usepackage[utf8]{inputenc}
\usepackage[T1]{fontenc}
\usepackage{lmodern}
\usepackage[portuguese]{babel}
\usepackage[a4paper,margin=2.5cm]{geometry}
\usepackage{xcolor}
\definecolor{TextEspresso}{HTML}{4A4238}
\definecolor{NudeTaupe}{HTML}{BFAFA6}
\definecolor{SoftBlush}{HTML}{D9C5B2}
\usepackage{titlesec}
\usepackage{enumitem}
\usepackage{listings}
\usepackage{hyperref}
\hypersetup{colorlinks=true, linkcolor=TextEspresso, urlcolor=TextEspresso}

\titleformat{\section}{\color{TextEspresso}\Large\bfseries}{\thesection}{0.5em}{}[{\color{NudeTaupe}\titlerule[0.8pt]}]
\titleformat{\subsection}{\color{TextEspresso}\large\bfseries}{\thesubsection}{0.5em}{}

\lstset{
  basicstyle=\ttfamily\small,
  breaklines=true,
  frame=single,
  rulecolor=\color{NudeTaupe},
  columns=fullflexible,
  showstringspaces=false
}

\newcommand{\guiaotitle}[2]{%
  \begin{center}
  {\Large\bfseries\color{TextEspresso} #1}\\[0.3cm]
  {\normalsize\color{NudeTaupe} #2}\\[0.2cm]
  {\color{NudeTaupe}\rule{4cm}{0.5pt}}
  \end{center}
  \vspace{0.5cm}
}

\begin{document}
\guiaotitle{Guião 11 -- NFS, iSCSI e Object Storage com MinIO}{Infraestruturas de Computação na Nuvem, Aula 11}

\section{Objetivos}
\begin{itemize}
    \setlength\itemsep{0.2cm}
    \item Observar os tipos de storage de rede suportados pelo Proxmox e a sua relação com a migração ao vivo.
    \item Configurar um servidor NFS numa VM e montar a partilha a partir de uma VM cliente.
    \item Configurar um alvo (target) iSCSI básico e ligar-se a partir de um initiator noutra VM.
    \item Demonstrar a diferença essencial entre acesso ao nível de ficheiro e ao nível de bloco, com acesso concorrente.
    \item Correr o MinIO e utilizar o cliente \texttt{mc} para criar um bucket e enviar/obter um objeto.
    \item Usar o NFS como PersistentVolume no cluster Kubernetes (opcional).
\end{itemize}

\section{Pré-requisitos e Material}
\begin{itemize}
    \setlength\itemsep{0.2cm}
    \item Duas VMs Linux no Proxmox: uma para servidor (\texttt{servidor}) e outra para cliente (\texttt{cliente}), na mesma rede. Pode reutilizar VMs de aulas anteriores.
    \item Acesso à interface web do Proxmox, para a Parte A.
    \item Docker instalado na VM onde se vai correr o MinIO (pode ser a mesma VM servidor).
    \item Acesso com privilégios de administração (\texttt{sudo}) em ambas as VMs.
    \item Opcionalmente, o cluster Kubernetes do Guião 05, para a Parte E.
    \item Pacotes necessários:
\begin{lstlisting}
# Na VM servidor
$ sudo apt update
$ sudo apt install -y nfs-kernel-server targetcli-fb

# Na VM cliente
$ sudo apt update
$ sudo apt install -y nfs-common open-iscsi
\end{lstlisting}
\end{itemize}

\section{Introdução}
Este guião cobre três formas distintas de armazenamento em rede, correspondentes aos conceitos vistos na teoria: NFS (ao nível de ficheiro, tipo NAS), iSCSI (ao nível de bloco, tipo SAN) e object storage (via MinIO, compatível com a API S3).

O objetivo não é apenas configurar cada um, mas perceber \emph{porque} diferem. O exercício mais importante do guião está na Parte C: montar a mesma partilha NFS em dois clientes ao mesmo tempo funciona; fazer o equivalente com iSCSI é uma receita para corromper dados. Compreender esta diferença é compreender a distinção entre ficheiro e bloco.

\section{Parte A -- Armazenamento de rede no Proxmox}
\begin{enumerate}[label=\arabic*.]
    \item \textbf{Observar os tipos de storage suportados.} Na interface web, aceda a \emph{Datacenter} $\rightarrow$ \emph{Storage} $\rightarrow$ \emph{Add} e consulte (sem adicionar nada) a lista de tipos disponíveis. Identifique os que correspondem a armazenamento de rede: \emph{NFS}, \emph{iSCSI}, \emph{CIFS/SMB}, \emph{Ceph RBD}, entre outros.

    \item \textbf{Relacionar com a migração ao vivo.} Consulte os storages atualmente configurados e a coluna que indica se são partilhados (\emph{Shared}):
\begin{lstlisting}
$ pvesm status
\end{lstlisting}
    Responda, para registo: se o disco de uma VM estiver num storage local, o que teria de acontecer para a migrar para outro nó? E se estivesse num storage partilhado por NFS ou Ceph? (Retome o conceito de migração ao vivo da Aula 03.)
\end{enumerate}

\section{Parte B -- Partilha de Ficheiros com NFS}
\begin{enumerate}[label=\arabic*.]
    \item Na VM \textbf{servidor}, criar a pasta a partilhar e definir permissões:
\begin{lstlisting}
$ sudo mkdir -p /srv/nfs/partilha
$ sudo chown nobody:nogroup /srv/nfs/partilha
\end{lstlisting}

    \item Editar \texttt{/etc/exports} para exportar a pasta para a rede do cliente (substituir pelo IP/sub-rede real):
\begin{lstlisting}
/srv/nfs/partilha 192.168.56.0/24(rw,sync,no_subtree_check)
\end{lstlisting}

    \item Aplicar a exportação e reiniciar o serviço:
\begin{lstlisting}
$ sudo exportfs -ra
$ sudo systemctl restart nfs-kernel-server
\end{lstlisting}

    \item Na VM \textbf{cliente}, criar o ponto de montagem e montar a partilha:
\begin{lstlisting}
$ sudo mkdir -p /mnt/nfs
$ sudo mount servidor:/srv/nfs/partilha /mnt/nfs
$ df -h /mnt/nfs
\end{lstlisting}

    \item Testar a escrita a partir do cliente e confirmar visibilidade no servidor:
\begin{lstlisting}
$ echo "ola via nfs" | sudo tee /mnt/nfs/teste.txt
# no servidor:
$ cat /srv/nfs/partilha/teste.txt
\end{lstlisting}

    \item \textbf{Testar o acesso concorrente.} Monte a \emph{mesma} partilha também no servidor (ou numa terceira VM, se disponível) e escreva a partir dos dois lados em simultâneo:
\begin{lstlisting}
# no cliente:
$ echo "escrito pelo cliente" | sudo tee /mnt/nfs/concorrente.txt
# no servidor (ou noutro cliente), imediatamente a seguir:
$ cat /srv/nfs/partilha/concorrente.txt
$ echo "acrescentado pelo servidor" | sudo tee -a /srv/nfs/partilha/concorrente.txt
# de volta ao cliente:
$ cat /mnt/nfs/concorrente.txt
\end{lstlisting}
    Confirme que ambos os lados veem as alterações um do outro. É o servidor NFS que gere o sistema de ficheiros e coordena os acessos -- guarde esta observação para comparar com o iSCSI.
\end{enumerate}

\section{Parte C -- Armazenamento em Bloco com iSCSI}
\begin{enumerate}[label=\arabic*.]
    \item Na VM \textbf{servidor}, criar um ficheiro de 500~MB para servir de disco virtual do target iSCSI:
\begin{lstlisting}
$ sudo mkdir -p /srv/iscsi
$ sudo truncate -s 500M /srv/iscsi/disco0.img
\end{lstlisting}

    \item Configurar o target iSCSI com \texttt{targetcli}:
\begin{lstlisting}
$ sudo targetcli
/> backstores/fileio create disco0 /srv/iscsi/disco0.img
/> iscsi/ create iqn.2026-11.icn.estga:servidor
/> iscsi/iqn.2026-11.icn.estga:servidor/tpg1/luns create /backstores/fileio/disco0
/> iscsi/iqn.2026-11.icn.estga:servidor/tpg1/acls create iqn.2026-11.icn.estga:cliente
/> saveconfig
/> exit
\end{lstlisting}

    \item Na VM \textbf{cliente}, definir o próprio IQN (deve corresponder ao ACL criado acima), em \texttt{/etc/iscsi/initiatorname.iscsi}:
\begin{lstlisting}
InitiatorName=iqn.2026-11.icn.estga:cliente
\end{lstlisting}

    \item Descobrir e ligar-se ao target a partir do cliente:
\begin{lstlisting}
$ sudo systemctl restart open-iscsi
$ sudo iscsiadm -m discovery -t sendtargets -p servidor
$ sudo iscsiadm -m node --targetname iqn.2026-11.icn.estga:servidor --login
\end{lstlisting}

    \item Confirmar que o novo disco de bloco apareceu e formatá-lo/montá-lo:
\begin{lstlisting}
$ lsblk
$ sudo mkfs.ext4 /dev/sdX   # substituir sdX pelo disco iSCSI identificado
$ sudo mkdir -p /mnt/iscsi
$ sudo mount /dev/sdX /mnt/iscsi
\end{lstlisting}

    \item \textbf{Observar a diferença essencial face ao NFS.} Repare no que acabou de fazer: teve de \emph{formatar} o dispositivo. No NFS isso nunca foi necessário, porque o sistema de ficheiros já existia do lado do servidor.

    Escreva um ficheiro e confirme:
\begin{lstlisting}
$ echo "ola via iscsi" | sudo tee /mnt/iscsi/teste.txt
$ ls -l /mnt/iscsi
\end{lstlisting}

    \textbf{Não faça isto:} montar este mesmo LUN em escrita a partir de um segundo cliente ao mesmo tempo. Cada cliente teria a sua própria cache e estruturas do sistema de ficheiros, sem saber da existência do outro, e o resultado seria a corrupção do ext4. É esta a diferença fundamental face ao NFS: no iSCSI o \emph{cliente} gere o sistema de ficheiros, e não há ninguém a coordenar acessos concorrentes.

    Para permitir escrita simultânea sobre um dispositivo de bloco partilhado seria necessário um sistema de ficheiros de cluster (ex.\ GFS2, OCFS2), com bloqueio distribuído.
\end{enumerate}

\section{Parte D -- Object Storage com MinIO}
\begin{enumerate}[label=\arabic*.]
    \item Correr o MinIO num container Docker:
\begin{lstlisting}
$ docker run -d --name minio -p 9000:9000 -p 9001:9001 \
    -e "MINIO_ROOT_USER=admin" -e "MINIO_ROOT_PASSWORD=minioadmin123" \
    -v /srv/minio-data:/data \
    minio/minio server /data --console-address ":9001"
\end{lstlisting}

    \item Instalar o cliente \texttt{mc} e configurar o alias para o servidor local:
\begin{lstlisting}
$ curl -O https://dl.min.io/client/mc/release/linux-amd64/mc
$ chmod +x mc
$ ./mc alias set localminio http://localhost:9000 admin minioadmin123
\end{lstlisting}

    \item Criar um bucket:
\begin{lstlisting}
$ ./mc mb localminio/dados-aula11
\end{lstlisting}

    \item Enviar (upload) um objeto para o bucket:
\begin{lstlisting}
$ echo "objeto de teste" > objeto.txt
$ ./mc cp objeto.txt localminio/dados-aula11/objeto.txt
\end{lstlisting}

    \item Listar o conteúdo do bucket e obter (download) o objeto de volta:
\begin{lstlisting}
$ ./mc ls localminio/dados-aula11
$ ./mc cp localminio/dados-aula11/objeto.txt objeto-descarregado.txt
$ cat objeto-descarregado.txt
\end{lstlisting}
\end{enumerate}

\section{Parte E -- NFS como PersistentVolume no Kubernetes (opcional)}
Esta parte assume o cluster do Guião 05 e reutiliza o servidor NFS configurado na Parte B.

\begin{enumerate}[label=\arabic*.]
    \item \textbf{Instalar o cliente NFS em todos os nós} do cluster:
\begin{lstlisting}
$ sudo apt install -y nfs-common
\end{lstlisting}

    \item \textbf{Criar o PersistentVolume e o PersistentVolumeClaim.} Crie \texttt{nfs-pv.yaml} (substitua o IP do servidor):
\begin{lstlisting}
apiVersion: v1
kind: PersistentVolume
metadata:
  name: nfs-pv
spec:
  capacity:
    storage: 1Gi
  accessModes:
    - ReadWriteMany
  nfs:
    server: <IP-do-servidor-NFS>
    path: /srv/nfs/partilha
  persistentVolumeReclaimPolicy: Retain
---
apiVersion: v1
kind: PersistentVolumeClaim
metadata:
  name: nfs-pvc
spec:
  accessModes:
    - ReadWriteMany
  resources:
    requests:
      storage: 1Gi
  storageClassName: ""
\end{lstlisting}
    Repare no modo de acesso \texttt{ReadWriteMany}: só é possível porque o NFS é armazenamento ao nível de ficheiro.

    \item \textbf{Aplicar e verificar a ligação (\emph{bound}):}
\begin{lstlisting}
$ kubectl apply -f nfs-pv.yaml
$ kubectl get pv,pvc
\end{lstlisting}

    \item \textbf{Montar o volume em dois Pods em nós diferentes} e confirmar que ambos veem os mesmos ficheiros:
\begin{lstlisting}
$ kubectl run escritor --image=busybox:1.36 --restart=Never \
    --overrides='{"spec":{"containers":[{"name":"c","image":"busybox:1.36","command":["sh","-c","echo escrito-pelo-pod > /dados/k8s.txt; sleep 3600"],"volumeMounts":[{"name":"v","mountPath":"/dados"}]}],"volumes":[{"name":"v","persistentVolumeClaim":{"claimName":"nfs-pvc"}}]}}'
$ kubectl run leitor --image=busybox:1.36 --restart=Never \
    --overrides='{"spec":{"containers":[{"name":"c","image":"busybox:1.36","command":["sh","-c","sleep 3600"],"volumeMounts":[{"name":"v","mountPath":"/dados"}]}],"volumes":[{"name":"v","persistentVolumeClaim":{"claimName":"nfs-pvc"}}]}}'
$ kubectl exec leitor -- cat /dados/k8s.txt
\end{lstlisting}

    \item \textbf{Confirmar no servidor NFS} que o ficheiro criado pelo Pod está lá:
\begin{lstlisting}
$ cat /srv/nfs/partilha/k8s.txt
\end{lstlisting}

    \item \textbf{Limpeza:}
\begin{lstlisting}
$ kubectl delete pod escritor leitor
$ kubectl delete -f nfs-pv.yaml
\end{lstlisting}
\end{enumerate}

\section{Entregável / Questões de Verificação}
\begin{enumerate}[label=\alph*)]
    \setlength\itemsep{0.2cm}
    \item Capturas de ecrã (ou output copiado) de: montagem NFS bem-sucedida no cliente, \texttt{lsblk} com o disco iSCSI montado, e \texttt{mc ls} a mostrar o objeto no bucket.
    \item No caso do NFS, o cliente acede aos dados ao nível de ficheiro ou de bloco? E no caso do iSCSI? Justifique com base no que foi montado em cada caso.
    \item Descreva o que observou no teste de acesso concorrente ao NFS (Parte B) e a razão pela qual o mesmo não deve ser feito com o volume iSCSI. Que componente, em cada caso, é responsável por gerir o sistema de ficheiros?
    \item Que vantagem tem o MinIO/S3 sobre o NFS quando se trata de armazenar milhões de ficheiros pequenos e não estruturados?
    \item O acesso ao MinIO foi feito através de que tipo de API? Compare com o modelo de acesso usado no NFS e no iSCSI.
    \item Se o disco de uma VM estiver num storage local do Proxmox, o que é necessário para a migrar para outro nó? E se estiver num storage partilhado? Relacione com a migração ao vivo (Aula 03).
    \item (Se fez a Parte E) Porque é que o PersistentVolume NFS pôde ser declarado com \texttt{ReadWriteMany}? Que modo de acesso teria de usar se o volume fosse iSCSI, e porquê?
    \item Indique, para cada um dos três mecanismos usados (NFS, iSCSI, MinIO), um cenário real de utilização em ambiente de cloud.
\end{enumerate}

\end{document}
